\documentclass[10pt,letterpaper]{article}
\usepackage{cornell-notes}

\title{Chapter 1: Foundations}
\author{}
\date{}
\setDocTitle{Chapter 1: Foundations}
\setDocSubtitle{Electronics}
\setDocOwner{Electronics Cornell Notes}
\setDocDate{\today}
\setCornellCollection{Electronics Cornell Notes}
\setCornellUnitType{Chapter}
\setCornellUnitNumber{1}
\setCornellUnitTitle{Foundations}
\hypersetup{pdftitle={Chapter 1: Foundations},pdfsubject={Cornell notes},pdfauthor={}}

\begin{document}
\maketitle

\begin{CornellOverviewBox}{Chapter Overview}
This chapter builds the working vocabulary of circuit design: voltage and current, resistive networks, signals, energy-storage elements, diodes, impedance, resonance, and practical passive hardware. The unifying habit is to replace a physical circuit with the simplest model that answers the current design question, then test the model against limits, loading, tolerances, and measurement conditions.

\textbf{Prerequisites.} Algebra, scientific notation, basic trigonometry, and comfort interpreting graphs versus time or frequency.

\textbf{Learning objectives.} Apply KCL, KVL, Ohm's law, and power relations; reduce linear networks with dividers and Thevenin equivalents; predict first-order transients; interpret impedance; and select passive components while accounting for loading, tolerance, power, and parasitics.
\end{CornellOverviewBox}

\begin{CornellNotesTable}
\CornellNoteRow{Design mindset}{Circuit analysis is not only equation solving. Begin with function, identify nodes and reference ground, estimate orders of magnitude, choose an ideal model, and then add the nonideal effect most likely to matter. Check conservation laws, limiting cases, component ratings, and physical plausibility.}
\CornellNoteRow{Voltage, current, and resistance}{Voltage is an energy-per-charge difference measured across two points; current is charge flow through a branch. KCL expresses charge conservation and KVL expresses energy conservation. For an ohmic resistor, $V=IR$ and $P=VI=I^2R=V^2/R$. A divider is accurate only when its load is included.}
\CornellNoteRow{Signals and levels}{Describe a signal by waveform, amplitude convention, offset, frequency, phase, source impedance, and reference node. For $v(t)=V_p\sin(2\pi ft+\phi)$, $V_{\mathrm{rms}}=V_p/\sqrt{2}$. Digital levels are voltage ranges with noise margins, not perfect binary numbers.}
\CornellNoteRow{Capacitors and RC circuits}{A capacitor obeys $i=C\,dv/dt$ and stores $E_C=CV^2/2$. In a first-order RC network, the distance from final value decays as $e^{-t/RC}$; $\tau=RC$ reaches about $63\%$ of a step. Real capacitors have leakage, ESR, ESL, dielectric absorption, tolerance, and polarity constraints.}
\CornellNoteRow{Inductors and transformers}{An inductor obeys $v=L\,di/dt$ and stores $E_L=LI^2/2$; its current cannot jump without impulse voltage. Real parts add winding resistance, core loss, saturation, parasitic capacitance, and coupling limits. Check isolation, insulation, core flux, frequency, and winding temperature.}
\CornellNoteRow{Diodes and rectification}{A diode conducts strongly in one direction after its exponential junction current becomes significant. Rectifiers convert polarity, reservoir capacitors trade ripple for pulsed charging current, and regulators reduce remaining variation. Check reverse voltage, average and surge current, recovery, leakage, capacitance, and heat.}
\CornellNoteRow{Impedance and reactance}{For sinusoidal steady state, $Z_R=R$, $Z_C=1/(j\omega C)$, and $Z_L=j\omega L$. Complex impedance preserves Ohm's-law form while encoding magnitude and phase. Resonance occurs when inductive and capacitive reactances cancel; resistance sets damping and $Q$.}
\CornellNoteRow{Practical passive parts}{Switches, relays, connectors, indicators, and variable components introduce contact resistance, bounce, isolation limits, lifecycle limits, polarity, pinout, and mating reliability. Select by worst-case electrical and environmental ratings rather than appearance.}
\CornellNoteRow{Markings and construction}{Confirm value, tolerance, power, package, polarity, orientation, and manufacturer pinout before assembly. Surface-mount parasitics, thermal handling, rework, creepage, and probe access are design considerations. Never infer identity from package shape alone.}
\end{CornellNotesTable}

\begin{CornellSummaryBox}{Chapter Synthesis}
Circuit foundations are a disciplined use of conservation laws and deliberately limited models. Resistive networks establish dc behavior; capacitors and inductors add state, time constants, phase, and stored energy; diodes add direction-dependent nonlinearity; and impedance unifies sinusoidal analysis. Practical success depends on source and load interaction, component limits, parasitics, and measurement loading as much as on nominal equations.
\end{CornellSummaryBox}

\begin{CornellExamBox}{Worked Examples}
\textbf{Loaded divider.} A $15\,\mathrm{V}$ source feeds $R_1=3.3\,\mathrm{k\Omega}$ over $R_2=6.8\,\mathrm{k\Omega}$ and a $20\,\mathrm{k\Omega}$ instrument. The lower leg is $R_2\parallel R_L=5.075\,\mathrm{k\Omega}$, giving $V_o=9.09\,\mathrm{V}$ instead of the unloaded $10.10\,\mathrm{V}$.

\textbf{RC step.} For $R=4.7\,\mathrm{k\Omega}$ and $C=2.2\,\mu\mathrm{F}$, $\tau=10.34\,\mathrm{ms}$. A step to $8\,\mathrm{V}$ reaches $8(1-e^{-2})=6.92\,\mathrm{V}$ at $2\tau$.
\end{CornellExamBox}

\begin{CornellWarningBox}{Safety and Limitations}
Discharge capacitors through a verified resistor before handling. Treat mains-referenced circuits, transformer primaries, high-current batteries, and inductive loads as hazardous. An oscilloscope ground clip is commonly tied to protective earth; connecting it to a live node can create a destructive short. Use isolation and differential measurement methods appropriate to the hazard.
\end{CornellWarningBox}

\begin{CornellExamBox}{Key Equation Sheet and Review}
\begin{itemize}
	\item $V=IR$; $P=VI=I^2R=V^2/R$.
	\item $i_C=C\,dv_C/dt$, $E_C=CV^2/2$, and $\tau_{RC}=R_{\mathrm{seen}}C$.
	\item $v_L=L\,di_L/dt$, $E_L=LI^2/2$, and $\tau_{RL}=L/R_{\mathrm{seen}}$.
	\item $Z_C=1/(j\omega C)$, $Z_L=j\omega L$, and $\omega_0=1/\sqrt{LC}$ for an ideal resonator.
	\item Verify reference node, source and load impedances, ratings, temperature, bandwidth, polarity, pinout, grounding, and discharge paths.
\end{itemize}
\end{CornellExamBox}

\end{document}
